\reportsection{Technical Overview}
\label{sec:attack}

% split into the a,b,c findings 
% maybe no lean yet here
% these are the considerations: ... this is how you calculate ...
% since we have these, we want a framework that allows them to have this
% begin the formalisation: concrete and machine checked. 

\paragraph{What the designer provides.}
A protocol built on \textsf{ark-vc} specifies four things: a field~$\field$,
a security level~$\secpar$, a hash function~$H$, and a choice of
interface, either a standard commitment or a hiding commitment.  These are the
only decisions the designer makes.  The security-requirements spine handles
everything derived from them: digest width, salt cardinality, parameter
certification, and the formal evidence that those parameters meet the stated
target.

\paragraph{What the spine returns.}
From $(\field,\secpar,H)$ and the interface choice, the parameter constructor
derives
\[
  D=\left\lceil\frac{\secpar+1+2h}{b}\right\rceil,\qquad
  k=\left\lceil\frac{\secpar}{b}\right\rceil,\qquad
  S=\left\lceil\frac{\secpar}{8}\right\rceil,
\]
where $b=\lfloor\log_2|\field|\rfloor$ is the conservative bits-per-field-element
and $h$ bounds the adversarial query budget.  These are not recommendations.
Each value is a checkable obligation: the constructor either certifies the
target is met or refuses the claim.  The designer cannot quietly under-specify;
the spine does not silently over-claim.

Field capacity enters the derivation directly, so the same $\secpar=128$ target
produces different outputs for different fields.  For BabyBear ($b=30$, $h=64$):
$D=9$, $k=5$, $S=16$.  For BN254 ($b\approx254$, $h=64$): $D=2$, $k=1$,
$S=16$.  These are not interchangeable.  The $44$-second binding break from
\cref{sec:motivation} used a $D=1$ configuration, one the constructor would
refuse.

\paragraph{The Lean~4 parallel.}
The Lean~4 formalisation gives the same derivation machine-checked standing.
For binding, the full chain (from an inconsistent accepted opening, through a
trace collision, to the birthday bound) is closed for any compliant
instantiation.  For hiding, the interface is typed to require explicit per-leaf
randomness; the cardinality prerequisites are proved; the remaining
distributional argument has a named, structured composition target.  A designer
consuming the formalisation does not re-derive or re-audit: the artefact is the
audit.  As new fields and hash functions enter \textsf{Arkworks}, the same
certificate framework applies without starting over.
